\section{Introduction}

Embodied-agent systems often fail at the boundary between planning and
execution. A planner can identify the right task-level intention yet still emit
an action name or argument structure that the executor cannot dispatch. Such
failures expose not only planning errors but also the planner-to-executor
contract, the action interface that mediates the handoff, and the runtime
validation policy that decides whether a bad first emission should terminate
the run or be repaired.

This execution boundary often remains implicit in embodied-agent studies, where
broader evaluations and stronger end-to-end models dominate the narrative
\cite{driess2023palme,brohan2023rt2,yang2025embodiedbench,lu2025isbench,kachaev2025mindmotionaligned}.
Those studies are valuable, but they entangle execution-stack design with
backend scale and benchmark breadth. We therefore ask a narrower
robotics-systems question: when the simulator setting and task families are
held fixed, how much do planner-to-executor contract design and lightweight
runtime validation affect invalid actions, recovery, success, and
planner/tool overhead?

We study two task families in Isaac Sim: language-conditioned navigation and
tabletop manipulation. The matrix compares three contract variants: an
under-specified direct-action contract (P0), a typed tool-call contract
(P1), and the same typed contract with a brief self-check (P2). It also
compares two runtime variants: bare dispatch (R0) and validation with a
single retry (R1). These variants are implemented through different planner
instructions, but the manipulated system variables are the action interface
presented to the executor and the runtime validation policy applied before
dispatch.

The findings are stable across the covered slices. P0/R0 is consistently the
weakest condition; adding runtime validation to that weak contract recovers
fragile invalid-first-action failures; explicit contracts P1 and P2
eliminate the invalid-action failure mode in the fixed-R0 slice; and P2 is
associated with lower planner/tool overhead than P1 in the covered
comparisons. Harder tasks increase planner and tool activity without reversing
the ordering. The paper is limited to this controlled simulator study and
draws bounded execution-design lessons rather than new-model or benchmark
claims.

This paper makes three contributions:

\begin{enumerate}
\item It presents a controlled Isaac Sim study that isolates two execution-design
axes, planner-to-executor contract design and lightweight runtime
validation, across navigation and tabletop manipulation tasks.
\item It provides evidence that typed tool-call contracts reduce invalid actions
relative to an under-specified direct-action contract, while runtime
validation with a single retry converts recoverable invalid-first actions
into successful executions in both covered task families.
\item It extracts a compact architectural lesson from the tested setup: contract
design, runtime validation, and a brief planner-side self-check are separate
systems levers that shape reliability and planner/tool overhead without
changing the underlying task matrix.
\end{enumerate}

Fig.~\ref{fig:system-overview} summarizes the evaluated boundary: a task
instruction reaches the planner, the planner emits an action request under one
of the three contract designs, the runtime applies either bare dispatch or
validate-and-retry, and the executor records the outcome metrics used
throughout the paper. The same figure also carries the compact failure-repair
trace that anchors the recovery discussion.

\begin{figure*}[!t]
  \centering
  \IfFileExists{\FigureAssetDir/fig1_system_overview_frozen.pdf}{
    \includegraphics[width=\textwidth]{\FigureAssetDir/fig1_system_overview_frozen.pdf}
  }{
    \includegraphics[width=\textwidth]{\FigureAssetDir/fig1_system_overview_frozen.png}
  }
  \caption{System overview of the evaluated planner-to-executor boundary. A task is mapped by the planner into one of three contract variants, then passed through bare dispatch \texttt{R0} or validate-and-retry \texttt{R1} before the executor (a lightweight deterministic backend for the main factorial, with a minimal Isaac Sim navigation slice for qualitative cross-backend checking). Insets show representative navigation and tabletop task scenes together with compact failure-repair traces in which non-dispatchable verbs are rejected and corrected into executor-visible tools.}
  \label{fig:system-overview}
\end{figure*}

